\documentclass[11pt]{article}
\usepackage[a4paper,margin=1in]{geometry}
\usepackage{amsmath,amssymb,amsthm,mathtools}
\usepackage{graphicx}
\usepackage{hyperref}
\usepackage{microtype}

\newtheorem{theorem}{Theorem}
\newtheorem{lemma}{Lemma}
\newcommand{\R}{\mathbb R}
\newcommand{\V}{\mathcal V}
\newcommand{\sech}{\operatorname{sech}}
\newcommand{\supp}{\operatorname{supp}}

\title{Exponential STFT concentration does not force a Poincar\'e inequality}
\author{Run \texttt{fa\_banach\_001}, agent lane 10}
\date{2026-08-11}

\begin{document}
\maketitle

\section*{Source and conjecture}

The source is Martin Rathmair, \emph{Stable STFT phase retrieval and
Poincar\'e inequalities} \cite{rathmair}.  It is available as
\href{https://arxiv.org/abs/2407.00398}{arXiv:2407.00398}.
Its short-time Fourier transform convention is
\[
 \V_gf(x,\xi)=\int_{\R}f(t)\overline{g(t-x)}e^{-2\pi i\xi t}\,dt.
\]
Conjecture 1.18 on page 7 fixes
\[
 g(t)=e^{t-e^t},\qquad
 \gamma(x,\xi)=e^{-a|x|-b|\xi|}\quad(a,b>0),
\]
and asks whether
\[
 \exists\delta>0:\quad |\V_gf(z)|e^{\delta|z|}\in L^\infty(\R^2)
 \quad\Longrightarrow\quad
 C_P\!\left((|\V_gf|^2*\gamma)^2\right)<\infty.
\]
Here \(C_P(w)\) is the least constant in
\[
 \inf_{c\in\mathbb C}\int_{\R^2}|u-c|^2w
 \le C_P(w)\int_{\R^2}|\nabla u|^2w,
 \qquad u\in C_b^\infty(\R^2).
\]

\begin{center}
\includegraphics[width=0.96\linewidth]{figures/open_problem_crop.png}
\end{center}

\section*{Proof intuition}

Strict positivity removes literally disconnected supports, but it does not
remove arbitrarily thin bottlenecks.  The ambiguity function
\(A=\V_gg\) is exponentially localized.  We place copies of \(g\) at the
lacunary positions \(R_n=L2^n\), with coefficients \(e^{-\kappa R_n}\).
Their exponential decay is enough to preserve the hypothesis of the
conjecture.  On the other hand, the smoothed spectrogram has a fixed-size
island near every \(R_n\), while its value halfway from \(R_{n-1}\) to \(R_n\)
has an additional exponentially small factor.  Step functions changing value
at those midpoints have non-negligible variance on the \(n\)-th island and
vanishingly smaller Dirichlet energy in the bottleneck.  Hence no uniform
Poincar\'e constant exists.

\section*{Counterexample}

Write \(T_Rh(t)=h(t-R)\) and \(A=\V_gg\).  The source computes
\begin{equation}\label{eq:Aexact}
 |A(x,\xi)|^2=\frac1{16}\sech^4(x/2)
 |\Gamma(2-2\pi i\xi)|^2.
\end{equation}
Put \(H(\xi)=|\Gamma(2-2\pi i\xi)|\).  In particular,
\begin{equation}\label{eq:Abound}
 |A(x,\xi)|\le C e^{-|x|}H(\xi),
 \qquad
 \sup_{\xi\in\R}e^{d|\xi|}H(\xi)<\infty\quad(0<d<\pi^2).
\end{equation}
The second assertion also follows immediately from the standard Gamma-function
identity used in the source.

\begin{theorem}\label{thm:counterexample}
For every \(a,b>0\), there is a nonzero \(f\in L^2(\R)\) such that
\[
 |\V_gf(z)|e^{\delta|z|}\in L^\infty(\R^2)
\]
for some \(\delta>0\), but, with
\(\gamma(x,\xi)=e^{-a|x|-b|\xi|}\) and
\(w=(|\V_gf|^2*\gamma)^2\), one has \(C_P(w)=\infty\).
Consequently Conjecture 1.18 is false for every parameter pair \(a,b>0\).
\end{theorem}

\begin{proof}
Fix \(r\in(0,2)\), \(s\in(0,\min\{a,r\})\), and
\begin{equation}\label{eq:kappachoice}
 0<\kappa<s/4.
\end{equation}
For a constant \(L>0\), to be chosen sufficiently large, set
\[
 R_n=L2^n,\qquad c_n=e^{-\kappa R_n},\qquad
 f=\sum_{n=1}^\infty c_nT_{R_n}g.
\]
Since \(\sum_nc_n<\infty\), this series converges in \(L^2(\R)\).
STFT covariance gives
\begin{equation}\label{eq:Fsum}
 F(x,\xi):=\V_gf(x,\xi)
 =\sum_{n=1}^\infty c_ne^{-2\pi i\xi R_n}A(x-R_n,\xi).
\end{equation}

\paragraph{The concentration hypothesis holds.}
By \eqref{eq:Abound},
\[
 |F(x,\xi)|\le CH(\xi)S(x),\qquad
 S(x):=\sum_nc_ne^{-|x-R_n|}.
\]
If \(0<\delta<\min\{\kappa,1,\pi^2\}\), then
\[
 e^{\delta|x|}S(x)
 \le\sum_nc_ne^{\delta R_n}e^{-(1-\delta)|x-R_n|}
 \le\sum_ne^{-(\kappa-\delta)R_n}<\infty.
\]
Together with \eqref{eq:Abound} and
\(e^{\delta|(x,\xi)|}\le e^{\delta|x|}e^{\delta|\xi|}\), this proves

\[
 |F(z)|e^{\delta|z|}\in L^\infty(\R^2).
\]

\paragraph{An upper bound in the bottlenecks.}
The centers \(R_n\) are uniformly separated.  Weighted Cauchy--Schwarz gives
\begin{align}
 S(x)^2
 &\le
 \left(\sum_nc_n^2e^{-r|x-R_n|}\right)
 \left(\sum_ne^{-(2-r)|x-R_n|}\right)\notag\\
 &\le C_r\sum_nc_n^2e^{-r|x-R_n|}.\label{eq:Ssquare}
\end{align}
Let
\[
 q=|F|^2*\gamma,\qquad w=q^2.
\]
Convolving \eqref{eq:Ssquare} separately in \(x\) and \(\xi\), and using
\[
 e^{-r|\cdot|}*e^{-a|\cdot|}\le C_{r,a,s}e^{-s|\cdot|},
\]
we obtain
\begin{equation}\label{eq:qbound}
 q(x,\xi)\le C B(\xi)P(x),\qquad
 B=H^2*e^{-b|\cdot|}\in L^2(\R),\quad
 P(x)=\sum_nc_n^2e^{-s|x-R_n|}.
\end{equation}
The \(L^2\)-claim for \(B\) follows from the exponential decay of \(H\) and
Young's inequality.

For \(n\ge2\), let
\[
 M_n=\frac{R_{n-1}+R_n}{2}=\frac32R_{n-1},\qquad
 I_n=[M_n-1,M_n+1].
\]
An elementary comparison of the terms in \(P\) shows that, uniformly for
\(x\in I_n\),
\begin{equation}\label{eq:midpoint}
 P(x)\le Cc_n^2e^{-\eta R_{n-1}},\qquad
 \eta:=s/2-2\kappa>0.
\end{equation}
Indeed, among \(j<n\) the largest term is \(j=n-1\), because
\(s-2\kappa>0\), and
\[
 \frac{c_{n-1}^2e^{-s(M_n-R_{n-1}-1)}}{c_n^2}
 =e^{-(s/2-2\kappa)R_{n-1}+s}.
\]
Among \(j\ge n\) the largest term is \(j=n\), and its ratio to \(c_n^2\)
is at most \(e^{-sR_{n-1}/2+s}\).  The remaining lacunary tails are bounded
geometric sums.  From \eqref{eq:qbound}--\eqref{eq:midpoint},
\begin{equation}\label{eq:gapmass}
 \int_{I_n\times\R}w(x,\xi)\,dx\,d\xi
 \le Cc_n^4e^{-2\eta R_{n-1}}.
\end{equation}

\paragraph{A lower bound on every island.}
Since \(A(0,0)\ne0\), choose a small rectangle
\(Q=[-\varepsilon,\varepsilon]^2\) and \(m>0\) such that
\(|A|\ge m\) on \(Q\).  Increasing \(L\), if necessary, makes the principal
term in \eqref{eq:Fsum} dominate all other translates, uniformly in \(n\):
\[
 |F(x,\xi)|\ge \frac m2c_n
 \quad\text{when }(x-R_n,\xi)\in Q.
\]
To see this, divide the upper bound for every other term by \(c_n\).  For
\(j<n\) it is, up to a fixed factor,
\(e^{-(1-\kappa)(R_n-R_j)}\), while for \(j>n\) it is
\(e^{-(1+\kappa)(R_j-R_n)}\).  Both lacunary sums tend uniformly to zero as
\(L\to\infty\).

Let \(Q'\) be a smaller concentric rectangle.  For
\(E_n=(R_n,0)+Q'\), integrate the convolution defining \(q\) over
\((R_n,0)+Q\).  The kernel \(\gamma\) has a fixed positive lower bound on the
relevant bounded difference set, and hence
\begin{equation}\label{eq:islandmass}
 q\ge d_0c_n^2\ \text{ on }E_n,\qquad
 \int_{E_n}w\ge d_1c_n^4
\end{equation}
with constants \(d_0,d_1>0\) independent of \(n\).

\paragraph{The Poincar\'e quotients diverge.}
The function \(q\) lies in \(L^1\cap L^\infty\): Moyal's identity gives
\(|F|^2\in L^1\), while \(\gamma\in L^1\cap L^\infty\).  Thus
\(w=q^2\in L^1\).  It is positive and nonzero; write
\(m_w=\int_{\R^2}w\in(0,\infty)\).

Choose \(\theta\in C^\infty(\R,[0,1])\) with
\(\theta(t)=0\) for \(t\le-1\), \(\theta(t)=1\) for \(t\ge1\), and bounded
derivative, and put \(u_n(x,\xi)=\theta(x-M_n)\).  By
\eqref{eq:gapmass},
\begin{equation}\label{eq:energy}
 \int_{\R^2}|\nabla u_n|^2w
 \le Cc_n^4e^{-2\eta R_{n-1}}.
\end{equation}
For every finite measure \(d\mu=w\,dz\),
\[
 \inf_c\int|u-c|^2\,d\mu
 =\frac1{2m_w}\iint|u(z)-u(z')|^2\,d\mu(z)d\mu(z').
\]
The mass of \(\{x\le M_n-1\}\) tends to \(m_w\), whereas \(u_n=1\) on
\(E_n\subset\{x\ge M_n+1\}\) for large \(n\).  Restricting the double integral
to these two sets and using \eqref{eq:islandmass} yields
\begin{equation}\label{eq:variance}
 \inf_c\int_{\R^2}|u_n-c|^2w\ge \frac{d_1}{2}c_n^4
\end{equation}
for all sufficiently large \(n\).  Dividing \eqref{eq:variance} by
\eqref{eq:energy} gives a quotient bounded below by
\(Ce^{2\eta R_{n-1}}\to\infty\).  Therefore \(C_P(w)=\infty\), completing
the counterexample.
\end{proof}

\section*{Verification notes}

The construction is analytic and uses no numerical or unproved dependency.
The parameter inequalities are simultaneously satisfiable for every
\(a,b>0\): choose \(r<2\), then \(s<\min\{a,r\}\), then
\(\kappa<s/4\), and finally \(L\) large.  The three review-critical estimates
are \eqref{eq:Ssquare}, the midpoint comparison \eqref{eq:midpoint}, and the
uniform main-translate dominance preceding \eqref{eq:islandmass}.

\section*{Novelty check and limitations}

On 2026-08-11, searches for the exact paper title, the phrases
\texttt{``Concentration implies finite Poincare constant''} and
\texttt{``exponential time-frequency concentration''}, and
\texttt{Conjecture 1.18} together with STFT/Poincar\'e terminology found the
source arXiv and published versions but no later claimed resolution.  The
local arXiv source corpus and the run's four lightweight indexes likewise had
no matching result.  This is a bounded novelty check, not a claim that every
non-arXiv source has been exhausted.

The theorem disproves the conjecture for all its parameter values.  It does
not characterize additional structural assumptions on \(f\) (for example,
unimodality or log-concavity of the spectrogram) that would restore a finite
Poincar\'e constant.

\section*{Human review recommendation}

Review as a likely-valid full counterexample.  The most useful independent
check is to verify the exponents in \eqref{eq:midpoint}: the previous island
contributes \(c_{n-1}^2e^{-s(R_n-R_{n-1})/2}\), while the target island has
mass scale \(c_n^2\); the choice \(\kappa<s/4\) makes their ratio decay.

\begin{thebibliography}{9}
\bibitem{rathmair}
Martin Rathmair,
\emph{Stable STFT phase retrieval and Poincar\'e inequalities},
arXiv:2407.00398; International Mathematics Research Notices 2024,
no.~22, 14094--14114.
\end{thebibliography}

\end{document}
